\paragraph{长度\trFiltration。}\label{sub:length-filtration}
设 $\scr{G}_0,\scr{G}_1$ 为 $\scr{Q}=\mathit{nu}\text{-}\mathit{fun}(\scr{A},\scr{B})$
的\trObject{}。所谓\textbf{长度\trFiltration}（length filtration），即
$\hom_{\scr{Q}}(\scr{G}_0,\scr{G}_1)$ 上的一个完备递减\trFiltration{} $F^{\bullet}$：
第 $F^r$ 项由那些满足 $T^0=\cdots=T^{r-1}=0$ 的预\trNaturalTransformation{} $T$
组成。令 $A=H(\scr{A})$、$B=H(\scr{B})$、$G_k=H(\scr{G}_k)$，则相伴的
\textbf{\trSpectralSequence}（spectral sequence）的首项为
\begin{equation}\label{eq:length-filtration-spectral-sequence}
E_1^{rs}=\prod_{X_0,\ldots,X_r}\Hom^{s}_{\mathbb{K}}\bigl(
  \Hom_{A}(X_{r-1},X_r)\otimes\cdots\otimes\Hom_{A}(X_0,X_1),
  \Hom_{B}(G_0X_0,G_1X_r)\bigr).
\end{equation}
尽管这个\trSpectralSequence{}一般未必\trConvergence，它仍可用于比较论证。在最初步的
层面上，\eqref{eq:length-filtration-spectral-sequence} 只涉及\trCohomology{}数据，
这一事实有如下推论：
\vkthmsync{6}
\begin{lemma}\label{thm:right-composition-iso}
设 $\scr{G}_0,\scr{G}_1,\scr{G}_2\in\Ob\scr{Q}$，并设 $T$ 是从 $\scr{G}_0$ 到
$\scr{G}_1$ 的一个\trNaturalTransformation{}。设对每个 $X\in\Ob\scr{A}$，用 $[T^0]$
作右\trComposition{}都给出\trIsomorphism{}%
$\Hom_{H(\scr{B})}(\scr{G}_1X,\scr{G}_2X)\to\Hom_{H(\scr{B})}(\scr{G}_0X,\scr{G}_2X)$。
那么，在 $H(\scr{Q})$ 中用 $T$ 作右\trComposition{}便是\trIsomorphism{}%
$\Hom_{H(\scr{Q})}(\scr{G}_1,\scr{G}_2)\to\Hom_{H(\scr{Q})}(\scr{G}_0,\scr{G}_2)$
（对左\trComposition{}有平行的结论）。\noqed
\end{lemma}
\begin{lemma}\label{thm:left-composition-cff}
若 $\scr{G}\colon\scr{A}\to\scr{B}$ 是\trCohomology\trFullyFaithful{}的\trNonunital{} \trAFunctor，则对任意
$\scr{C}$，$\scr{L}_{\scr{G}}\colon\mathit{nu}\text{-}\mathit{fun}(\scr{C},\scr{A})
\to\mathit{nu}\text{-}\mathit{fun}(\scr{C},\scr{B})$ 亦然（关于
$\scr{R}_{\scr{G}}\colon\mathit{nu}\text{-}\mathit{fun}(\scr{B},\scr{C})
\to\mathit{nu}\text{-}\mathit{fun}(\scr{A},\scr{C})$ 的相应命题一般不真，但当
$\scr{G}$ 为\trQuasiisomorphism{}时成立）。\noqed
\end{lemma}
由\eqref{eq:natural-transformation-boundary} 可见，作用在
\eqref{eq:length-filtration-spectral-sequence} 上的\trDifferential{}%
$\partial=\partial^{rs}$ 为
\begin{equation*}
\begin{split}
(\partial T)(a_{r+1},\ldots,a_1)
&=(-1)^{(r+s)(|a_1|-1)+1}\,T(a_{r+1},\ldots,a_2)G_0(a_1)\\
&\quad+(-1)^{\clubsuit_r+(r+s)}\,G_1(a_{r+1})T(a_r,\ldots,a_1)\\
&\quad+\sum_{n}(-1)^{\clubsuit_{n+1}+(r+s-1)}\,
T(a_{r+1},\ldots,a_{n+2}a_{n+1},\ldots,a_1).
\end{split}
\end{equation*}
（这里的规定与我们惯常的略有不同：$T$ 是 $E_1^{rs}$ 的元素，因而在\trCohomology{}%
层面上起作用；类似地，诸 $a_i$ 是 $A$ 中的\trMorphism{}。）我们将把这个
双\trGraded\trComplex{}记作 $\mathit{CC}^{r+s}(A,B)^s=E_1^{rs}$，把它的
$\partial$-\trCohomology{}记作 $\mathit{HH}^{r+s}(A,B)^s=E_2^{rs}$，并分别称之为
$A$ 以 $B$ 为\trCoefficients{}的\textbf{\trHochschild{}\trComplex}（Hochschild complex）
与\textbf{\trHochschildCohomology}（Hochschild cohomology）（这依赖于诸\trFunctor{}%
$G_0$、$G_1$，尽管记号本身并未显示这一点）。第一个非平凡的列是
\begin{equation*}
E_2^{0s}=\mathit{HH}^s(A,B)^s=\Hom^{s}_{\mathit{Nu}\text{-}\mathit{Fun}(A,B)}(G_0,G_1),
\end{equation*}
而\textbf{边\trHomomorphism}（edge homomorphism）
$H^s(\hom_{\scr{Q}}(\scr{G}_0,\scr{G}_1))\to E_2^{0s}$ 即 $T\mapsto H(T)$。
于是，其余各列中的\trHochschildCohomology{}群可视为\eqref{eq:H-of-nu-fun-to-Nu-Fun}
是否\trFullyFaithful{}的\trObstruction。
\begin{remark}\label{thm:hochschild-classical-definition}
\trHochschildCohomology{}的经典定义是对一个 $\mathbb{K}$-代数 $A$ 连同其上一个\trBimodule{}%
$M$ 而言的，它给出一个\trGraded{}群 $HH(A,M)$。这可以按直接的方式
推广到\trGraded{}代数与\trBimodule，此时 \trHochschildCohomology{}便成为
双\trGraded{}的。作为特殊情形，设有两个代数\trHomomorphism{} $G_0,G_1\colon A\to B$。
于是 $B$ 成为 $A$-\trBimodule（经由 $G_1$ 从左侧作用，经 $G_0$ 从右侧作用），
从而有\trHochschildCohomology{}群 $HH(A,B)$。前文所考虑的群，正是把这个概念从代数
（可视为只有一个\trObject{}的\trCategory{}）推广到一般\trCategory{}的显然扩张，
只是我们对双\trGraded{}块的标记方式多少有些非常规（我们将从更抽象的\trHomology{}代数的
观点，在下文注记 2.7 中继续这一讨论）。

为避免术语上的混淆，我们指出：在 \trAinfty{}世界中还存在另外的\trHochschildCohomology{}%
理论，即对一个 \trAinfty{}代数连同其上的一个\trAinfty\trBimodule{}%
所配的\trGraded{}群 $HH(\scr{A},\scr{M})$。同样地，若
$\scr{G}_0,\scr{G}_1\colon\scr{A}\to\scr{B}$ 是 \trAinfty{}代数之间的\trHomomorphism，
便可用它们把 $\scr{B}$ 变成 $\scr{A}$ 上的\trAinfty\trBimodule，从而定义
$HH(\scr{A},\scr{B})$。而把它推广到\trACategory，得到的恰是\trCohomology{}%
$H(\hom_{\scr{Q}}(\scr{G}_0,\scr{G}_1))$——它正是我们的\trSpectralSequence{}的
目标。
\end{remark}
